\section{Two metrics and controlled pole support}\label{g:analytic}\label{rb:two-metrics}

The section problem involves two line bundles with different metric
requirements.  The first supplies a finite volume measure and controls
cotangent tensors.  The second is the bundle in which we seek a section;
its bounded metric is preserved when the problem is transferred to a
lower-dimensional variety.

All varieties in this section are smooth connected projective varieties
 over $\mathbb C$, unless a normal intermediate model is specified.
We use additive notation for line bundles.  A semipositive singular
Hermitian metric has local squared norm $e^{-\phi}$, where $\phi$ is
plurisubharmonic and not identically $-\infty$.  Bounded weights are
locally bounded from both sides.  Let $D\geq0$ be an integral divisor and
put $A=-K_V+D$.  Its canonical section $s_D$ is an $A$-valued top form.
For a metric $h_A$ on $A$, the notation
\[
                 \mu=|s_D|_{h_A}^{2}
\]
denotes the associated positive measure.  The finite-volume condition
$\int_V\mu<\infty$ is a separate assumption; it does not follow from
semipositivity of $h_A$.  This measure has positive finite density almost
everywhere and gives no mass to proper algebraic subsets.

\subsection{Models and adjoint measures}\label{g:models}

For a dominant rational map with connected general fibres we use a
normal projective intermediate model $\bar V$ and a smooth resolution
$W$ fitting into
\[
 W\longrightarrow\bar V,\qquad
 \bar p:\bar V\longrightarrow V,\quad
 \bar f:\bar V\longrightarrow Y,
 \qquad p:W\longrightarrow V,\quad f:W\longrightarrow Y.
\]
Here $Y$ is smooth projective, $\bar p$ is birational, and $\bar f$ is
surjective and equidimensional.  Such models follow by applying
flattening to the graph, resolving the modified base, taking its
dominating transform, and normalizing \cite{RaynaudGruson1971}.
The dimension bound from the flat model persists under base change for
the dominating transform and under finite normalization.  Every
vertical prime of $\bar V$ therefore maps onto a prime divisor of $Y$.
A prime of $W$ mapping into a subset of $Y$ of codimension at least two
is exceptional over $\bar V$ and hence over $V$.  The smooth resolution
itself need not be equidimensional.

Write
\[
 J=K_{W/V}\geq0,\qquad S=J+p^*D.
\]
The section $s_S=s_Jp^*s_D$ is a section of $K_W+p^*A$.  Its squared
norm is the birational pullback of $\mu$, including the squared Jacobian.
Change of variables off the exceptional locus gives a finite measure
on $W$.  This statement concerns densities almost everywhere and does
not require a pointwise definition of a singular metric on its polar
set.

\subsection{Transferring the two metrics}

The next statement separates normalization of divisor orders from the
analytic input which gives curvature.  This permits both a rank-one
adjoint argument and a direct calculation along torus orbits.

\begin{lemma}[Transfer with controlled pole support]\label{g:transfer}\label{g:transfer-recalled}\label{rb:transfer}
Let $V$ be smooth connected projective, let $D\geq0$ be integral,
and suppose $A=-K_V+D$ has a semipositive singular metric $h_A$ with
$\int_V|s_D|_{h_A}^2<\infty$. Let $L$ have a semipositive metric $h_L$
with bounded weights.  Use the model diagram above with
$0<\dim Y\leq\dim V$.  Suppose a nonzero rational section $\rho$ of $kL$,
$k>0$, has no horizontal poles on $\bar V$.  On a dense base open set
put
\[
 g_0(y)\,d\lambda_y=f_*p^*\mu,
 \qquad
 S_\rho(y)=\operatorname*{ess\,sup}_{W_y}|p^*\rho|_{h_L^k}^2.
\]
Assume $-\log g_0$ in coordinate frames and $-\log S_\rho$ have psh
representatives on this open set.  Then, after clearing a positive
integer denominator, there are a line bundle $B$ with bounded
semipositive weights, an effective integral divisor $D_Y$, and a
regular nonzero map
\[
                   f^*B\longrightarrow bk\,p^*L
\]
for some $b>0$.  The line bundle $-K_Y+D_Y$ has a semipositive metric
whose distinguished adjoint measure has finite mass.  Every nonzero section
of $aB+cD_Y$, where $a>0$ and $c\geq0$ are integers, produces a
nonzero rational section of $abkL$ with poles only on $D$.
If the data and $\rho$ are invariant for a torus acting trivially on
$Y$, then the produced section is invariant.
\end{lemma}

\begin{proof}
We first normalize the section, obtaining both the regular map and its
bounded base metric. We then correct the pushforward measure along base
primes whose poles disappear on return to $V$ or remain within $D$.

For a prime $T$ of $Y$ set
\[
 d_T=\min_{Q\mapsto T}
       \frac{\operatorname{ord}_Q(\bar p^*\rho)}
            {\operatorname{ord}_Q(\bar f^*T)},
 \qquad B_0=\sum_Td_TT.
\]
The orders in this expression are taken on $\bar V$.  Only finitely
many $d_T$ are nonzero.  Equidimensionality and the absence of
horizontal poles give
\[
             \operatorname{div}_{\bar V}(\rho)-\bar f^*B_0\geq0.
\]
This is an effective $\mathbb Q$-Cartier divisor.  Choose $b>0$ clearing
the coefficients of $B_0$ and set $B=bB_0$.  Normality extends the
normalized section in codimension one to a regular morphism
$\bar f^*B\to bk\bar p^*L$, and pullback gives the displayed map.
Over each prime $T$ a component realizing the minimum is absent from
its vanishing divisor.

For a local meromorphic equation $l$ of $B$, define the weight
\[
                     -b\log S_\rho+\log|l|^2.
\]
It is psh on the initial open set.  Its negative exponential is the
fibre supremum of the squared norm of the regular normalized section
$(p^*\rho)^b/f^*l$.  Properness and bounded weights of $h_L$ bound this
norm above locally over every base point, giving a lower bound for the
weight even near the omitted locus. At a general point of each
omitted base prime $T$ choose a minimizing component. The normalized section
is nonzero there. At general points of this component and $T$, the
transverse map has coordinates $y_1=z_1^e$, $y_i=z_i$ for
$2\leq i\leq\dim Y$. Thus every nearby fibre meets a patch where
the normalized section norm is bounded below. The fibre
measure has positive density almost everywhere on such a patch, so
the essential supremum has a fixed positive lower bound.  This gives
an upper bound for the weight near a general point of that prime.
Psh removable singularities extend the weight there and then across
the remaining analytic set of codimension at least two.  The lower
bound persists, while the extended psh weight is locally bounded
above.  These weights glue and give the bounded metric on $B$.

The bounded metric is now constructed. For the finite-volume metric,
call a base prime $T$ invisible if every component of $\bar f^*T$ is
exceptional over $V$ or maps into $D$.  There are finitely many such
primes, since they are images of the finitely many exceptional or
boundary primes. Choose a component above $T$ on $\bar V$ and denote
its strict transform on $W$ by $Q$. At its general point write
\[
 e=\operatorname{ord}_Q(f^*T),\qquad
 a=\operatorname{ord}_Q(J+p^*D).
\]
The local psh coefficient weight of $h_A$ is bounded above, so its
negative exponential is bounded below by a positive constant. The
Jacobian and boundary section therefore bound the pulled-back measure
below by a positive constant times $|z_1|^{2a}$ times coordinate volume.
Integrating over
a fixed disc in the free variables and using $y_1=z_1^e$ gives
\begin{equation}\label{rb:density-lower}
              g_0(y)\geq C|y_1|^{2(a-e+1)/e}
\end{equation}
almost everywhere near the general base point.  If $T$ is visible,
choose a component with $a=0$; then $-\log g_0$ is bounded above there.
For the invisible primes choose sufficiently large integral
coefficients of $D_Y$.  Equation~\eqref{rb:density-lower} makes
\[
                  -\log g_0+\log|s_{D_Y}|^2
\]
locally bounded above at their general points.  It extends psh in
codimension one and then across codimension two, defining a metric on
$-K_Y+D_Y$.  Pairing it with $s_{D_Y}$ cancels the added factor.  Its
adjoint measure is $g_0d\lambda_y=f_*p^*\mu$ almost everywhere and is
therefore finite.

A section of $aB+cD_Y$ is a rational section of $aB$ with poles only on
$D_Y$.  Pull it back and compose with the $a$-th power of the normalized
map.  On $\bar V$ all possible new poles lie above invisible primes.
They are thus exceptional over $V$ or map into $D$.  The normal
intermediate model excludes nonexceptional divisors over base subsets
of codimension at least two.  Birational pushdown therefore leaves
poles only on $D$, and a sufficiently large multiple of $D$ clears
them.  The rational section remains nonzero.  The normalization uses
only base factors, so it preserves invariance when present.
\end{proof}

\begin{lemma}[Rank-one sufficient condition]\label{g:moments}\label{rb:rank-one}
In the setting of Lemma~\ref{rb:transfer}, suppose on the generic fibre
of $f$ that
\[
 h^0\bigl(W_\eta,(J+p^*D+mkp^*L)|_{W_\eta}\bigr)=1
                       \qquad(m\geq0).
\]
Then the two psh assumptions in that lemma hold.  There is also an
invariant version: take $D=0$, smooth input metrics invariant under a
compact form of an acting torus, and an invariant $\rho$; it suffices
that the invariant dimensions in the display equal one.
\end{lemma}

\begin{proof}
On a base open set where $\rho$ is regular, define
\[
 g_m(y)\,d\lambda_y
       =f_*\bigl(|p^*\rho|_{h_L^k}^{2m}p^*\mu\bigr).
\]
These are finite positive densities almost everywhere, since the
section norms are bounded on compact preimages and $p^*\mu$ has finite
mass. Give $p^*(A+mkL)$ the coefficient metric
$h_m=p^*(h_Ah_L^{mk})$, and consider the integrable direct image
\[
 \mathcal F_m=f_*\bigl(\mathcal O_W(K_{W/Y}+p^*(A+mkL))
                         \otimes\mathcal I(h_m)\bigr).
\]
For a local canonical frame $\eta_Y$ on the base, the section
\[
             u_m=\frac{s_{J+p^*D}(p^*\rho)^m}{f^*\eta_Y}
\]
is locally integrable on the total space and therefore belongs to
$\mathcal F_m$. It is nonzero and generates
the ordinary direct image generically by the rank-one hypothesis.
Consequently $\mathcal F_m$ also has rank one. On a dense open set
where $u_m$ is its frame, the squared direct-image norm is $g_m$.

Singular direct-image positivity gives $-\log g_m$ a psh
representative there. We use the formulation with the actual
multiplier ideal in \cite[Theorem~21.1]{HPS}, based on
\cite[Theorem~3.3.5]{PT2018}. Equivalently, in the relative Bergman
description, the weight is the logarithm of the generator's squared
holomorphic coefficient minus $\log g_m$ pulled back from the base
\cite[Theorem~0.1]{BerndtssonPaun2008}.
In the invariant case, use smooth direct-image positivity
\cite[Theorem~1.2]{Berndtsson2009}.  Averaging over the compact torus gives
a holomorphic projection onto the invariant line.  Distinct weight
spaces are orthogonal, so its subspace metric equals the quotient
metric and is semipositive.

There is a fixed dense open set where $f$ is a submersion and each
fibre has a patch avoiding the zeros of $s_{J+p^*D}$ and $\rho$.
Integrating over such a patch gives
$g_m\geq Cc^m$ locally, with $C,c>0$.  Hence the psh representatives of
$-m^{-1}\log g_m$ have uniform local upper bounds, independent of
possible extra base-change exceptions for individual $m$; psh
extension fills those exceptions on the fixed open set.  Finite
fibre measure and its positive density almost everywhere imply
\[
  -\frac1m\log g_m(y)\longrightarrow
  -\log\operatorname*{ess\,sup}_{W_y}|p^*\rho|_{h_L^k}^2
\]
almost everywhere. The limit is finite there by the bounded metric
and the regularity of $\rho$. Local compactness for psh functions
now gives the required representative. Indeed, a subsequence tending
locally uniformly to $-\infty$ would contradict this finite limit.
Every $L^1_{\mathrm{loc}}$-convergent subsequence has a further
subsequence converging almost everywhere, so its psh limit agrees
almost everywhere with the displayed power-mean limit. The direct-image
argument at $m=0$ supplies $-\log g_0$.
\end{proof}


\subsection{Divisors from twisted forms}

\begin{lemma}\label{g:determinant}\label{rb:determinant}
Let $L$ have bounded semipositive weights on a smooth projective
$n$-fold $V$ with $H^0(V,\Omega_V^j)=0$ for $j>0$.  There are a line
bundle $M$ mapping nontrivially into a tensor power of $\Omega_V^1$ and
effective integral divisors
\[
                       N_i\sim M+m_iL,
                       \qquad m_i\longrightarrow+\infty.
\]
The integers $m_i$ can be chosen strictly increasing and divisible by any prescribed positive integer.
\end{lemma}

\begin{proof}
Replacing $L$ at the outset by a prescribed positive multiple gives the divisibility assertion; it is therefore enough to prove the statement without that prescription. Hodge symmetry gives $\chi(\mathcal O_V)=1$, so Serre duality gives
$\chi(K_V)=(-1)^n$.  Riemann--Roch makes $\chi(K_V+mL)$ a polynomial
in $m$ with nonzero constant term.  It is nonzero for all sufficiently
large integers.  A fixed cohomological degree $q$ therefore has
$H^q(V,K_V+mL)\ne0$ for unbounded $m>0$.  The metric on $mL$ has trivial
multiplier ideal.  Hard Lefschetz with pseudoeffective coefficients
\cite[Theorem~0.1]{DPS2001} gives nonzero sections of
$\Omega_V^{n-q}\otimes mL$ at these exponents.

Take the saturated subsheaf $\mathcal E\subset\Omega_V^{n-q}$ defined
by the span of their untwisted generic images.  Its determinant
$M=(\bigwedge^{\operatorname{rank}\mathcal E}\mathcal E)^{**}$ is a line
bundle and maps to a cotangent tensor power.  This map is first defined
outside codimension two and extends because the ambient tensor bundle
is locally free.  Every nonzero vector from the family can be completed
to a generic basis using vectors of the same family.  Wedging that
basis gives a nonzero section of $M$ twisted by the sum of its positive
exponents.  These sums are unbounded because one basis vector can
have an arbitrarily large exponent.  Saturation ensures that the
maps factor through $\mathcal E$, and extension across codimension two
gives global determinant sections.  Their zero divisors are the
required $N_i$, after taking a strictly increasing subsequence.  This
is the determinant method of \cite[Lemma~4.1]{LP2018}, used for
anticanonical nonvanishing in \cite[Lemma~5.1]{LMPTX2023}.
\end{proof}
